\setcounter{page}{1}
\pagenumbering{roman}
\begin{center}
\textbf{\large{Lời cam đoan}	}
\end{center}
\addcontentsline{toc}{chapter}{Lời cam đoan}
Tôi là Trần Anh Đức, sinh viên lớp K67I-IT2. Tôi xin cam đoan khoá luận tốt nghiệp này là do tôi tự nghiên cứu, cài đặt với sự nỗ lực của bản thân và đặc biệt là dưới sự hướng dẫn của TS. Lê Khánh Trình. Tôi cam kết tất cả những tài liệu tham khảo đã được liệt kê ở phần cuối khoá luận là nội dung tham khảo nằm trong giới hạn, phạm vi cho phép theo quy chế của Trường Đại học Công nghệ - Đại học Quốc gia Hà Nội. Nếu không đúng sự thật, tôi xin hoàn toàn chịu trách nhiệm.

\begin{flushright}
Hà Nội, ngày 5 tháng 5 năm 2026
\end{flushright}

\begin{changemargin}{12cm}{2cm}
Sinh viên
\\[2cm]
\end{changemargin}

\begin{changemargin}{11.5cm}{2cm}
Trần Anh Đức
\end{changemargin}

\newpage\cleardoublepage

\begin{center}
\textbf{\large{Tuyên bố về sự hỗ trợ của AI}}
\end{center}
\addcontentsline{toc}{chapter}{Tuyên bố về sự hỗ trợ của AI}

Trong quá trình thực hiện khoá luận, tôi có sử dụng Claude (Anthropic) như một công cụ hỗ trợ cho các công việc sau:

\begin{itemize}
    \item Nghiên cứu các bridge blockchain, khảo sát các mô hình và giải pháp hiện có, từ đó tổng hợp ưu nhược điểm để làm cơ sở cho việc thiết kế sản phẩm.
    \item Hỗ trợ cài đặt theo yêu cầu và gỡ lỗi, với mục tiêu giảm thiểu thời gian viết syntax và tập trung vào phần logic nghiệp vụ (business logic) chính của hệ thống.
\end{itemize}

Toàn bộ ý tưởng, thiết kế kiến trúc, thuật toán cốt lõi cũng như nội dung trình bày trong khoá luận đều do tôi tự xây dựng và chịu trách nhiệm. Claude chỉ đóng vai trò công cụ hỗ trợ, không thay thế cho công việc nghiên cứu và phát triển.

\begin{flushright}
Hà Nội, ngày 5 tháng 5 năm 2026
\end{flushright}

\begin{changemargin}{12cm}{2cm}
Sinh viên
\\[2cm]
\end{changemargin}

\begin{changemargin}{11.5cm}{2cm}
Trần Anh Đức
\end{changemargin}